In den Dokumentklassen~\texttt{article} und \texttt{report} wird nur die
Fußzeile genutzt. Dort steht mittig zentriert die Seitennummer. Die
Dokumentklasse~\texttt{book} verwendet hingegen lebende Kolumnentitel.
Dabei erscheint in der Kopfzeilen neben der Seitenzahl auch der aktuelle
Gliederungspunkt. Mittels des Pakets~\ctanref{fancyhdr} kann dieses Verhalten
auch für andere Klassen aktiviert werden. In der Datei~\texttt{latex/fancyhdr.tex}
der Vorlage ist bereits alles dafür vorbereitet. Sie müssen nur die
Kommentarzeichen vor der Passage in Listing~\ref{lst:final:fancyhdr}
entfernen.

\lstinputlisting[language={[LaTeX]{TeX}},style=block,float,caption={Aktivierung lebender Kolumnentitel},label={lst:final:fancyhdr}]{code/final_fancyhdr.tex}

Die Einstellungen gelten für das gesamte Dokument. Auf der ersten Seite jedes
Kapitels wird dennoch die einfache Form der Seitennummierung angewendet.
Sollten Sie auch den Vorspann der Arbeit in der einfachen Variante nummerieren
wollen, so schalten Sie mit dem Kommando \verb+\pagestyle+ auf den Seitenstil
\texttt{plain} um, wie es in Listing~\ref{lst:final:fancyhdr2} gezeigt wird.

\lstinputlisting[language={[LaTeX]{TeX}},style=block,float,caption={Abweichende Nummierierung für den Vorspann},label={lst:final:fancyhdr2}]{code/final_fancyhdr2.de.tex}
